% Section 2: Observers, fibres, and the residual principle
\label{sec:observers-body}

Everything in this paper is relative to what some observer is allowed to see.  This
section fixes that relativity once, so that later statements about "what cannot be
recovered" have a definite scope.  The material is developed on one deliberately
small worked model, the binary affine monoid; nothing in \S\ref{sec:triadic-inference}
or later depends on that model being the only one available.

\subsection{The observer contract}

\begin{definition}[Observer contract]\label{def:observer-contract}
Let $\cD$ be a class of processes.  A \emph{finite observer contract} on $\cD$ is a
tuple
\[
  Q=(F_Q,\equiv_Q,\Pi_Q,B_Q)
\]
consisting of
\begin{enumerate}
\item an \emph{observation map} $F_Q\colon\cD\to\cO_Q$ into a codomain of readable
  data, possibly many-sorted;
\item an \emph{equivalence} $\equiv_Q$ on $\cO_Q$, declaring when two readable data
  are indistinguishable to $Q$;
\item a family of \emph{projectors} $\Pi_Q$, each specifying which part of an
  observation is retained by a given reading;
\item a \emph{budget} $B_Q$, bounding the finite computations $Q$ may perform.
\end{enumerate}
\end{definition}

The budget is part of the contract, not an implementation detail.  A claim of the
form "no observer of class $\cO$ can determine $X$" has no content unless $\cO$ is
fixed, and $B_Q$ is where the finiteness of the observation enters.

\begin{remark}
The four components are not independent.  Enlarging a budget can refine
$\equiv_Q$, and adding a projector can shrink $F_Q$'s image.  Every statement below
that quantifies over observers therefore names the contract it uses.
\end{remark}

\subsection{Fibres and characteristics}

\begin{definition}[Observation fibre]\label{def:observation-fibre}
For $o\in\cO_Q$ and a reading $D\in\{t,X,\cK\}$ (temporal, spatial, constructive), the
\emph{fibre} of $o$ under $D$ is
\[
  \Fib_{D,Q}(o)=\{\,p\in\cD : F_{D,Q}(p)\equiv_Q o\,\}.
\]
\end{definition}

\begin{definition}[Characteristic]\label{def:characteristic}
A \emph{characteristic} of $p$ for the contract $Q$ is a triple
\[
  \chi_Q(p)=\bigl(\bar\chi_Q(p),\,\pi_p,\,R_{Q}(p)\bigr),
\]
where $\bar\chi_Q(p)$ is the operational quotient (the readable value),
$\pi_p$ is a certificate that $p$ maps to it, and $R_Q(p)$ is the \emph{residual}:
the data retained about $p$ that the quotient does not determine.
\end{definition}

\begin{remark}[A characteristic is not a value]
Definition~\ref{def:characteristic} is where this paper departs from the usual
usage.  A spectrum, a gradient, an eigenvalue, or a scalar is at best
$\bar\chi_Q$; the residual is part of the object, and propositions below exhibit
cases where discarding it destroys exactly the information a later reading needs.
\end{remark}

\subsection{The residual principle}

The following principle is used as an admissibility test throughout; it is the
Part~I counterpart of the rule that no quotient may be used without a declared
residual.

\begin{proposal}[Residual principle]\label{prop:residual-principle}
A quotient $\bar\chi_Q$ of a class $\cD$ may be used as a state for a declared class
of continuations and observations only when accompanied by
\begin{enumerate}
\item an exact reconstruction statement for the continuations in question, or
\item a residual $R_Q$ whose contents are characterized, together with a stated
  sense in which the continuation cannot depend on what $R_Q$ omits.
\end{enumerate}
\end{proposal}

This is a \textsf{Proposal} rather than a theorem because options (i) and (ii) are
requirements on presentation, not mathematical assertions.  Its bite is in what it
forbids: a quotient presented without either is \emph{not} eligible as a state.

\subsection{Identifiability is relative to a hypothesis class}

The worked model is the binary affine monoid.  Let
\[
  f_{k,b}(x)=2^{k}x+b,\qquad k,b\in\Z,
\]
with composition in execution order.  Writing $S$ for $f_{0,1}$ (successor) and $D$
for $f_{1,0}$ (doubling), every element has a normal form $D^{k};S^{b}$.

\begin{lemma}[Affine product law]\label{lem:affine-product}
For $f_{i}=f_{k_i,b_i}$ executed in the order $f_1$ then $f_2$,
\[
  f_2\odot f_1
  =\bigl(k_1+k_2,\ 2^{k_2}b_1+b_2\bigr),
\]
with identity $f_{0,0}$.  In particular
\[
  \chi(S;D)=(1,2)\neq(1,1)=\chi(D;S).
\]
\end{lemma}

\begin{proof}
$f_2(f_1(x))=2^{k_2}(2^{k_1}x+b_1)+b_2=2^{k_1+k_2}x+2^{k_2}b_1+b_2$.  The two
composites differ in their second coordinate because $2\cdot1+0=2$ while
$2^{0}\cdot 0+1=1$.
\end{proof}

Identifiability of $(k,b)$ is not an absolute property; it is relative to a declared
class of hypotheses and to the probes the contract allows.

\begin{proposition}[Two-probe temporal recovery]\label{prop:two-probe}
Probing the temporal reading at $x=0$ and $x=1$ recovers $(k,b)$ uniquely within the
affine class: with $y_i=f_{k,b}(i)$ one has $b=y_0$ and $2^{k}=y_1-y_0$.
\end{proposition}

\begin{proof}
Immediate from $f_{k,b}(0)=b$ and $f_{k,b}(1)=2^{k}+b$.
\end{proof}

\begin{remark}[One probe is not enough]\label{rem:one-probe}
A single probe at $x=1$ cannot separate $S$ from $D$, since $f_{0,1}(1)=2=f_{1,0}(1)$.
The two-probe result is therefore a statement about the contract, not about the
process.
\end{remark}

The same phenomenon appears on the spatial side, and the contrast is instructive:
there a single \emph{exact interval} suffices, whereas a radius alone does not.

\begin{proposition}[Spatial recovery from one exact interval]\label{prop:interval-recovery}
Let the spatial reading act on closed intervals by inverse image,
$\rho_X(k,b)\bigl(I(m,r)\bigr)=I\bigl((m-b)/2^{k},\,r/2^{k}\bigr)$.  From the single
observation $I(m,r)\mapsto I(m',r')$ one recovers $2^{k}=r/r'$ and hence
$b=m-2^{k}m'$.
\end{proposition}

\begin{proof}
The observed image has midpoint $(m-b)/2^{k}$ and radius $r/2^{k}$.  Dividing gives
$2^{k}$, and substituting gives $b$.
\end{proof}

\begin{warning}
Proposition~\ref{prop:interval-recovery} does not say that radii determine
translations.  With the radius alone the observation is $r\mapsto r/2^{k}$, which is
blind to $b$ entirely.  The informative observer is doing the work, not the
mathematics of affine maps.
\end{warning}

\subsection{The four-valued outcome, and what ``unknown'' means}

Inference in this setting is not a function into answers; it is a computation that
may fail to conclude.  We therefore insist on four outcomes.

\begin{definition}[Outcome discipline]\label{def:outcome-discipline}
An inference query returns exactly one of \emph{unique}, \emph{ambiguous},
\emph{inconsistent}, or \emph{unknown}.  An empty candidate set is reported as
\emph{inconsistent} only when exhaustion of the declared search space is certified;
a budget expiry or an incomplete enumeration is reported as \emph{unknown}.
\end{definition}

\begin{warning}[Failure to find is not a proof of absence]
\label{warn:unknown-not-nonexistence}
Definition~\ref{def:outcome-discipline} forbids the inference
``search found nothing, therefore nothing exists''.  This prohibition recurs in
Part~II, where an exhausted search and a certified empty fibre are deliberately
distinguished (\S\ref{sec:adequacy}).
\end{warning}

\subsection{The quotient is lossy, and the loss is the point}

It is tempting to treat a characteristic as a compressed description of a process.
The affine model shows immediately that compression and loss are the same phenomenon.

\begin{proposition}[Quotient without section]\label{prop:lossy-quotient}
Let $\chi_\cK$ be the constructive characteristic that records the parameter pair of
a normal form.  Then
\[
  \chi_\cK(S;D)=\chi_\cK(D;S;S)=(1,2),
\]
so $\chi_\cK$ does not distinguish the raw histories, although
$\chi_\cK\circ\operatorname{NF}=\id$ holds on parameter pairs.  Consequently
$\operatorname{NF}\circ\chi_\cK\neq\id$ on raw histories.
\end{proposition}

\begin{proof}
The composites are $f_{1,0}\circ f_{0,1}=f_{1,2}$ and
$f_{0,1}\circ f_{0,1}\circ f_{1,0}=f_{1,2}$ by Lemma~\ref{lem:affine-product}, while
the words $S;D$ and $D;S;S$ are distinct.
\end{proof}

\begin{proposition}[Compression is inseparable from loss]\label{prop:compression-loss}
The characteristic of $S;D^{n}$ is the pair $(n,2^{n})$, whose canonical word
$D^{n};S^{2^{n}}$ has $2^{n}$ successor symbols, while the parameter pair requires
$O(n)$ bits.  The compression ratio is therefore exponential, and it is purchased
exactly by the identification in Proposition~\ref{prop:lossy-quotient}.
\end{proposition}

\begin{proof}
By induction on $n$ using Lemma~\ref{lem:affine-product}: $f_{n,2^{n}}\odot S=f_{n,2^{n}+1}$
and doubling increments the first coordinate.  The normal form of $(n,2^{n})$ is
$D^{n};S^{2^{n}}$ by construction.
\end{proof}

\begin{remark}
Propositions~\ref{prop:lossy-quotient} and~\ref{prop:compression-loss} are the
smallest instance of the pattern that organises the whole paper: a quotient is
useful in proportion to what it forgets, and admissible in proportion to whether the
forgotten part is accounted for.
\end{remark}

\subsection{Summary of what this section fixes}

The section has introduced the observer contract, the fibre, the characteristic with
its residual, and the outcome discipline.  Two things should be carried forward.

First, \emph{identifiability is a property of a contract}, not of a process
(Proposition~\ref{prop:two-probe} and Remark~\ref{rem:one-probe}).

Second, \emph{a quotient must be presented with a residual}
(Proposition~\ref{prop:residual-principle} and
Proposition~\ref{prop:lossy-quotient}).  Every later section that computes a
characteristic states which of the two options in
Proposition~\ref{prop:residual-principle} it uses.
